% Introdução curta, compartilhada entre a aula 1 e o roteiro de bancada.
\begingroup
\begin{frame}{Nossa tecnologia na prática: gpdk045}
\begin{center}\begin{tikzpicture}[x=1cm,y=1cm,font=\small]
\draw[blue_theme,thick,fill=blue_theme!6] (3.7,0) rectangle (6.4,2.2);
\node[blue_theme,font=\bfseries] at (5.05,1.35){gpdk045};\node[font=\scriptsize] at (5.05,.65){CMOS genérico · 45 nm};
\node[align=center,font=\scriptsize] (a) at (1.1,3){\textbf{Schematic}\\Símbolos e terminais};
\node[align=center,font=\scriptsize] (b) at (9,3){\textbf{Simulação}\\Modelos elétricos};
\node[align=center,font=\scriptsize] (c) at (1.1,-.55){\textbf{Layout}\\Células parametrizadas};
\node[align=center,font=\scriptsize] (d) at (9,-.55){\textbf{Verificação}\\Regras DRC/LVS};
\draw[seta] (3.7,1.7)--(a);\draw[seta] (6.4,1.7)--(b);\draw[seta] (3.7,.3)--(c);\draw[seta] (6.4,.3)--(d);
\path[use as bounding box] (-.5,-1.15) rectangle (10.6,3.65);
\end{tikzpicture}\end{center}
{\small\centering Na nossa montagem: $V_{DD}=1$ V, $W=120$ nm, $L=45$ nm.\par
\texttt{nmos1v} / \texttt{pmos1v} · um finger · multiplicador 1.\par}
\fonte{PDK = Process Design Kit. Fonte: \href{https://community.cadence.com/cadence_technology_forums/f/custom-ic-design/36796/basics-of-pdk/1349734}{Cadence — Basics of PDK}: kit genérico para ensino e demonstração. Parâmetros aplicados: Quick Start da prática; 1 V é a alimentação usada neste minicurso.}
\end{frame}

\begin{frame}{O mesmo inversor CMOS, três geometrias}
\tresInversoresGeometria
\fonte{Vista superior conceitual: PMOS em verde, NMOS em azul, Gate em vermelho e contatos/metal em cinza. As setas W/L indicam o canal. Os valores valem para ambos os transistores em cada exemplo, com um finger e multiplicador 1. Contatos, poços e espaçamentos simplificados; não é um layout fabricável. As duas variações são propostas de comparação, sem novas curvas simuladas.}
\end{frame}

\begin{frame}{Como W/L aparece na corrente de Drain?}
\small\centering
\textcolor{blue_theme}{\textbf{NMOS em saturação · aproximação de canal longo}}\par\medskip
\[
 I_{DS}\approx\frac{1}{2}\,\mu_n C_{\mathrm{ox}}\,
 \textcolor{blue_theme}{\frac{W}{L}}\,(V_{GS}-V_{TN})^2
\]
{\scriptsize $V_{GS}>V_{TN}$ e $V_{DS}\geq V_{GS}-V_{TN}$; corpo ligado ao Source.}\par\medskip
$\mu_n$: mobilidade\qquad $C_{\mathrm{ox}}$: capacitância por área\par
$V_{TN}$: limiar\qquad $W/L$: razão geométrica\par\medskip
\begin{tabular}{lccc}\toprule
Modelo ideal & Referência & Dobrar W & Dobrar L\\\midrule
Corrente relativa & $1\times$ & $2\times$ & $0{,}5\times$\\\bottomrule
\end{tabular}
\fonte{Mesmas tensões, mobilidade, Cox e limiar. Para o módulo da corrente PMOS, usar $\mu_p$, $V_{SG}$ e $|V_{TP}|$. Fonte: \href{https://web.mit.edu/6.101/www/s2020/handouts/L06.pdf}{MIT 6.101, aula 6, pp. 19–20}. A lei quadrática explica tendências; a simulação em 45 nm usa os modelos do PDK e não precisa seguir exatamente essas proporções. W maior também aumenta capacitâncias.}
\end{frame}
\endgroup
